% ============================================================
\chapter*{Pr\'eface}
\addcontentsline{toc}{chapter}{Pr\'eface}
% ============================================================

\section*{Objectifs et motivation}

La th\'eorie de la mesure et de l'int\'egration constitue l'un des piliers fondamentaux
de l'analyse math\'ematique moderne. N\'ee des travaux pionniers d'\'Emile Borel et
d'Henri Lebesgue au d\'ebut du XX\textsuperscript{e} si\`ecle, cette th\'eorie a
profond\'ement transform\'e notre compr\'ehension des concepts de longueur, d'aire, de
volume, et plus g\'en\'eralement de \og taille \fg d'un ensemble. L'int\'egrale de
Lebesgue, qui en d\'ecoule, \'etend consid\'erablement l'int\'egrale de Riemann et
fournit le cadre naturel pour la th\'eorie des probabilit\'es, l'analyse fonctionnelle,
les \'equations aux d\'eriv\'ees partielles et de nombreuses autres branches des
math\'ematiques.

Ce cours s'adresse aux \'etudiants de niveau Master (M1/M2) en math\'ematiques
fondamentales ou appliqu\'ees. Il suppose une ma\^itrise solide de l'analyse r\'eelle
(suites, s\'eries, continuit\'e, d\'erivabilit\'e, int\'egrale de Riemann) ainsi que des
notions de base de topologie g\'en\'erale. Une familiarit\'e avec l'alg\`ebre lin\'eaire
et les espaces vectoriels norm\'es est \'egalement souhaitable.

\section*{Probl\'ematique historique}

L'int\'egrale de Riemann, d\'evelopp\'ee dans les ann\'ees 1850, pr\'esente des limitations
fondamentales. Consid\'erons la fonction de Dirichlet~:
\[
f(x) = \mathbf{1}_{\Q}(x) = \begin{cases} 1 & \text{si } x \in \Q, \\ 0 & \text{si } x \notin \Q. \end{cases}
\]
Cette fonction n'est pas Riemann-int\'egrable, alors qu'intuitivement, puisque $\Q$
est \og n\'egligeable \fg dans $\R$, on s'attendrait \`a ce que $\int_0^1 f = 0$.

De m\^eme, si $(f_n)$ est une suite de fonctions Riemann-int\'egrables convergeant
simplement vers une fonction $f$, il n'est en g\'en\'eral pas vrai que
\[
\lim_{n \to \infty} \int_a^b f_n(x)\, dx = \int_a^b f(x)\, dx,
\]
m\^eme lorsque le membre de droite a un sens. Pour que l'\'echange limite-int\'egrale
soit valide au sens de Riemann, on a besoin de la convergence \emph{uniforme}, ce
qui est une condition tr\`es restrictive.

L'int\'egrale de Lebesgue r\'esout ces deux probl\`emes de mani\`ere \'el\'egante~: la
fonction de Dirichlet devient int\'egrable (d'int\'egrale nulle), et le th\'eor\`eme de
convergence domin\'ee permet l'\'echange limite-int\'egrale sous des hypoth\`eses bien
plus faibles que la convergence uniforme.

\section*{Id\'ees directrices}

L'approche de Lebesgue repose sur une id\'ee fondamentalement diff\'erente de celle
de Riemann. Alors que Riemann d\'ecoupe le \emph{domaine} de la fonction en petits
intervalles, Lebesgue d\'ecoupe l'\emph{image} de la fonction et mesure les
pr\'eimages~:

\begin{intuition}[title=\textbf{Philosophie de Lebesgue}]
Imaginez que vous devez compter une grande somme d'argent compos\'ee de pi\`eces de
diff\'erentes valeurs. La m\'ethode de Riemann consiste \`a prendre les pi\`eces une par
une, dans l'ordre o\`u elles se trouvent. La m\'ethode de Lebesgue consiste \`a trier
d'abord les pi\`eces par valeur, puis \`a compter combien il y en a de chaque type.
Le r\'esultat est le m\^eme, mais la seconde m\'ethode est plus souple et s'applique
\`a davantage de situations.
\end{intuition}

Pour que cette approche fonctionne, il faut pouvoir \og mesurer \fg les
pr\'eimages $f^{-1}([a, b))$. Cela n\'ecessite~:
\begin{enumerate}
    \item Une th\'eorie des \textbf{ensembles mesurables} (les $\sigma$-alg\`ebres),
    qui d\'etermine quels ensembles peuvent \^etre mesur\'es.
    \item Une th\'eorie de la \textbf{mesure}, qui attribue une \og taille \fg
    \`a chacun de ces ensembles.
    \item Une d\'efinition des \textbf{fonctions mesurables}, dont les pr\'eimages
    d'ensembles raisonnables sont mesurables.
    \item La construction de l'\textbf{int\'egrale} par approximation via des
    fonctions \'etag\'ees (\emph{fonctions simples}).
\end{enumerate}

\section*{Plan du cours}

Le cours est organis\'e en douze chapitres qui suivent cette progression~:

\begin{description}
    \item[Chapitre 1.] $\sigma$-Alg\`ebres et espaces mesurables~: les structures
    fondamentales. Tribus engendr\'ees, tribu bor\'elienne.
    \item[Chapitre 2.] Mesures~: d\'efinitions, exemples, $\sigma$-additivit\'e,
    mesures finies et $\sigma$-finies, lemme de continuit\'e.
    \item[Chapitre 3.] Th\'eor\`eme de Carath\'eodory~: construction de mesures via
    des mesures ext\'erieures, passage d'une pr\'e-mesure \`a une mesure compl\`ete.
    \item[Chapitre 4.] Mesure de Lebesgue sur $\R^n$~: construction,
    propri\'et\'es d'invariance, ensembles non mesurables.
    \item[Chapitre 5.] Fonctions mesurables~: d\'efinitions, stabilit\'e par
    op\'erations alg\'ebriques et limites, approximation par fonctions simples.
    \item[Chapitre 6.] Int\'egrale de Lebesgue~: construction en trois \'etapes
    (fonctions simples, fonctions positives, fonctions int\'egrables).
    \item[Chapitre 7.] Th\'eor\`emes de convergence~: convergence monotone (Beppo
    Levi), lemme de Fatou, convergence domin\'ee (Lebesgue).
    \item[Chapitre 8.] Espaces $L^p$~: in\'egalit\'es de H\"older et Minkowski,
    compl\'etude, s\'eparabilit\'e, dualit\'e.
    \item[Chapitre 9.] Mesures sign\'ees~: d\'ecomposition de Hahn et de Jordan,
    variation totale.
    \item[Chapitre 10.] Th\'eor\`eme de Radon--Nikodym~: absolue continuit\'e,
    d\'eriv\'ee de Radon--Nikodym, d\'ecomposition de Lebesgue.
    \item[Chapitre 11.] Mesures produits~: $\sigma$-alg\`ebre produit,
    th\'eor\`emes de Fubini et Tonelli.
    \item[Chapitre 12.] Mesures de Radon et th\'eor\`eme de repr\'esentation de Riesz.
\end{description}

\section*{Conventions et notations}

\begin{itemize}
    \item $\N = \{0, 1, 2, \ldots\}$ d\'esigne l'ensemble des entiers naturels
    (incluant z\'ero).
    \item $\R^+ = [0, +\infty)$ et $\overline{\R}^+ = [0, +\infty]$.
    \item La droite r\'eelle \'etendue est $\overline{\R} = [-\infty, +\infty]$.
    \item $\mathbf{1}_A$ d\'esigne la fonction indicatrice de $A$.
    \item $f^+ = \max(f, 0)$, $f^- = \max(-f, 0)$, $f = f^+ - f^-$,
    $\abs{f} = f^+ + f^-$.
    \item $A_n \uparrow A$ signifie $(A_n)$ croissante et $A = \bigcup_n A_n$.
    \item $A_n \downarrow A$ signifie $(A_n)$ d\'ecroissante et $A = \bigcap_n A_n$.
    \item \og p.p.\fg signifie \og presque partout \fg (en dehors d'un ensemble
    de mesure nulle).
    \item $(X, \mathcal{A}, \mu)$ d\'esigne un espace mesur\'e.
    \item $\mathcal{B}(\R^n)$ d\'esigne la tribu bor\'elienne de $\R^n$.
    \item $\lambda^n$ d\'esigne la mesure de Lebesgue sur $\R^n$ ($\lambda = \lambda^1$).
\end{itemize}

\section*{Pr\'erequis}

\begin{itemize}
    \item \textbf{Analyse r\'eelle}~: propri\'et\'es de $\R$, suites et s\'eries
    num\'eriques, continuit\'e, d\'erivabilit\'e, int\'egrale de Riemann, s\'eries de
    fonctions, convergence uniforme.
    \item \textbf{Topologie g\'en\'erale}~: espaces m\'etriques, ouverts, ferm\'es,
    compacit\'e, connexit\'e.
    \item \textbf{Alg\`ebre lin\'eaire}~: espaces vectoriels, normes, espaces de
    Banach (notions de base).
    \item \textbf{Th\'eorie des ensembles}~: op\'erations ensemblistes, axiome du
    choix (pour la construction d'ensembles non mesurables).
\end{itemize}

\section*{R\'ef\'erences bibliographiques}

\begin{itemize}
    \item H.~L.~Royden et P.~M.~Fitzpatrick, \emph{Real Analysis}, 4\textsuperscript{e}
    \'ed., Pearson, 2010.
    \item W.~Rudin, \emph{Real and Complex Analysis}, 3\textsuperscript{e} \'ed.,
    McGraw-Hill, 1987.
    \item G.~B.~Folland, \emph{Real Analysis: Modern Techniques and Their Applications},
    2\textsuperscript{e} \'ed., Wiley, 1999.
    \item P.~Halmos, \emph{Measure Theory}, Springer, 1974.
    \item H.~Br\'ezis, \emph{Analyse fonctionnelle}, Masson, 1983.
    \item J.~Neveu, \emph{Bases math\'ematiques du calcul des probabilit\'es},
    Masson, 1970.
    \item C.~Zuily et H.~Queff\'elec, \emph{Analyse pour l'agr\'egation}, Dunod, 2013.
    \item R.~L.~Schilling, \emph{Measures, Integrals and Martingales}, Cambridge
    University Press, 2005.
\end{itemize}

\bigskip
\hfill \emph{L'auteur}

\hfill \emph{Mars 2026}
